\chapter{Guía de Despliegue y Manual de Usuario}
\label{anx:despliegue}

Este anexo describe cómo levantar \textbf{Seshat} y operarlo de principio a fin. El despliegue de
referencia se empaqueta como un conjunto de contenedores autocontenido, de modo que un único comando
arranca la base de datos, los servicios auxiliares y la aplicación. La última sección recorre la
interfaz web como manual de usuario.

\section{Requisitos previos}
\label{anx:despliegue-requisitos}

Para minimizar las dependencias que recaen sobre quien lo despliega, \textbf{Seshat} se distribuye
completamente dockerizado. Por eso, el único requisito para arrancar el sistema completo es
\textbf{Docker} con el complemento \textbf{Compose}~v2 (\texttt{docker compose}); los contenedores
aportan todo lo demás —PostgreSQL, el almacén de secretos y de \textit{blobs}, el servidor de trazas
y los dos servicios de Seshat—, sin instalar nada más en la máquina anfitriona. El único estado que
sobrevive a los contenedores vive en volúmenes de Docker, de modo que detenerlos y volver a
levantarlos no pierde ni la base de datos ni las trazas. La sección~\ref{anx:despliegue-arranque}
detalla cada servicio y el orden en que arrancan.

Para el desarrollo —ejecutar los tests, los arneses de evaluación o la \acs{API} fuera de los
contenedores— se requiere además \textbf{Python~3.13} o superior y el gestor de paquetes
\texttt{uv}\footnote{%
  Gestor de paquetes y proyectos de Python escrito en Rust: \url{https://docs.astral.sh/uv/}.%
}. A partir del fichero \texttt{uv.lock} versionado, \texttt{uv} reconstruye el entorno de forma
determinista: instala las versiones exactas allí fijadas y resuelve hasta el modelo de spaCy que
sostiene las heurísticas de puntuación de confianza (sección~\ref{sec:confidence-scoring}), que de
otro modo exigiría una descarga aparte.

\section{Configuración de credenciales}
\label{anx:despliegue-credenciales}

La configuración se suministra mediante dos ficheros de entorno que se copian de sus plantillas
versionadas:

\begin{itemize}
  \item \texttt{.env} (a partir de \texttt{.env.example}) cumple dos funciones: Docker Compose lo
        lee automáticamente para sustituir las variables de infraestructura del
        \texttt{docker-compose.yml} —credenciales de PostgreSQL, puertos, nombre del \textit{bucket},
        región—, y \texttt{python-dotenv} lo carga en las ejecuciones locales fuera de los
        contenedores, donde añade las claves de \acs{API} y los ajustes que apuntan los servicios a
        \texttt{localhost}.
  \item \texttt{.env.docker} (a partir de \texttt{docker/.env.docker.example}) se inyecta
        explícitamente en los contenedores \texttt{seshat-api} y \texttt{localstack}. Contiene las
        claves de \ac{API} —que se siembran en el almacén de secretos al arrancar— y los ajustes de
        configuración de Seshat que sobrescriben los valores por defecto.
\end{itemize}

Las credenciales no se leen directamente del entorno en tiempo de ejecución, sino a través de un
\textbf{resolutor de secretos} seleccionable por configuración. Con el proveedor \texttt{env} se
leen de las variables de entorno, adecuado para ejecuciones locales; con el proveedor \texttt{aws}
se recuperan de AWS Secrets Manager. El despliegue en contenedores usa este último apuntando a
\textbf{LocalStack}, que emula el servicio: así el arranque local ejercita exactamente el mismo
camino de código que un despliegue en la nube, sin depender de infraestructura real.

Las claves mínimas para un arranque funcional son la clave raíz que gobierna la administración
(\texttt{SESHAT\_ROOT\_API\_KEY}) y las de los proveedores del pipeline por defecto:
\texttt{OPENAI\_API\_KEY} (\textit{embeddings} y \textit{grounding}), \texttt{ANTHROPIC\_API\_KEY}
(identificación y resolución) y \texttt{ASSEMBLYAI\_API\_KEY} (transcripción).

\section{Arranque con Docker Compose}
\label{anx:despliegue-arranque}

Una vez copiados y rellenados los dos ficheros de entorno, el conjunto completo de contenedores se
levanta desde la raíz del proyecto con un único comando:

\begin{center}
  \ttfamily docker compose up
\end{center}

En la primera ejecución, Docker Compose construye las imágenes de Seshat que aún no existen\footnote{%
  Compose sólo construye las imágenes ausentes; tras modificar el código fuente hay que forzar la
  reconstrucción con \texttt{docker compose up -{}-build}, o construir las imágenes manualmente
  con \texttt{docker compose build <service\_name>}. Alternativamente, se puede crear un fichero
  \texttt{docker-compose.override.yml} y definir en él los ajustes personalizados, tales como
  volúmenes a las carpetas donde se encuentra el código fuente. %
} y arranca los seis servicios de la tabla~\ref{tab:compose-services} en el orden que imponen sus
dependencias. El servicio
\texttt{migrate} aplica las migraciones de Alembic (\texttt{alembic upgrade head}), que crean los
esquemas \texttt{ops} y \texttt{knowledge\_base} descritos en la sección~\ref{anx:db-schema}, y
termina; \texttt{seshat-api} sólo arranca una vez que \texttt{migrate} ha concluido con éxito y que
PostgreSQL, LocalStack y MLflow están sanos, y \texttt{seshat-ui} espera a su vez a que la \ac{API}
lo esté.

\begin{table}[H]
  \centering
  \small
  \begin{tabularx}{\textwidth}{@{}l l c >{\raggedright\arraybackslash}X@{}}
    \toprule
    \textbf{Servicio} & \textbf{Origen} & \textbf{Puerto} & \textbf{Función} \\
    \midrule
    \texttt{postgres}   & \texttt{pgvector/pgvector:pg17} & 5432 & Base de datos: esquemas \texttt{ops} y \texttt{knowledge\_base} más los \textit{embeddings} de pgvector. \\
    \texttt{migrate}    & Dockerfile (\texttt{migrate})   & —    & Aplica las migraciones de Alembic una sola vez y termina. \\
    \texttt{mlflow}     & \texttt{mlflow:v3.12.0}         & 5000 & Servidor de trazas y métricas de observabilidad. \\
    \texttt{localstack} & \texttt{localstack:3}           & 4566 & Emula AWS Secrets Manager y S3; siembra los secretos y crea el \textit{bucket} al arrancar. \\
    \texttt{seshat-api} & Dockerfile (\texttt{api})       & 8000 & \ac{API} REST y cola de trabajos en proceso. \\
    \texttt{seshat-ui}  & Dockerfile.ui                   & 8501 & Interfaz web (Streamlit) sobre la \ac{API}. \\
    \bottomrule
  \end{tabularx}
  \caption{Servicios de \texttt{docker compose} y el puerto que cada uno expone en la
           máquina anfitriona.}
  \label{tab:compose-services}
\end{table}

El arranque de la \ac{API} atraviesa además dos comprobaciones que actúan como salvaguardas: un
sondeo a todos los proveedores de modelos configurados y la \textbf{verificación de la puerta de
evaluación}, que aborta el arranque si no existe un \texttt{eval\_gate.json} con veredicto favorable
(véase el capítulo~\ref{ch:evaluacion}).

En el despliegue en local ambas comprobaciones se desactivan mediante las variables de entorno
\texttt{API\_\_SKIP\_EXTERNAL\_PROVIDER\_PING} y \texttt{API\_\_SKIP\_EVAL\_GATE}, de modo
que la demostración arranca sin exigir una ejecución previa de los arneses ni que todos los
proveedores estén disponibles. En un despliegue real ambas se dejan activas: se ejecuta antes
\texttt{seshat eval} para producir una puerta favorable y sólo entonces se permite servir tráfico.

Una comprobación adicional, ajena al \textit{gate} pero también resuelta por configuración, es la
verificación TLS de las llamadas salientes a proveedores externos. En una red corporativa cuyo proxy
inyecta un certificado raíz propio, ese certificado suele estar provisionado en el almacén de
confianza del sistema operativo pero no en el que Python empaqueta por defecto, lo que hace fallar
la verificación aun cuando el tráfico es legítimo. La variable de entorno \texttt{USE\_OS\_TRUSTSTORE}
resuelve esto sin desactivar la verificación: activada, hace que el módulo \texttt{ssl} de Python
delegue en el almacén de confianza del sistema operativo en lugar del suyo propio, de modo que las
llamadas a través de ese proxy quedan verificadas igual que las demás.

Una vez sanos todos los servicios, la interfaz web queda accesible en \texttt{localhost:8501}, la
\ac{API} en \texttt{localhost:8000} y la interfaz de MLflow en \texttt{localhost:5000}. Fuera de los
contenedores, los mismos dos pasos se ejecutan a mano con la línea de comandos: \texttt{seshat
migrate} aplica las migraciones y \texttt{seshat api} levanta el servidor.

\section{Manual de usuario}
\label{anx:despliegue-manual}

Los usuarios operan \textbf{Seshat} a través de una interfaz web: el cliente fino de Streamlit cuyo
diseño se motiva en la sección~\ref{sec:ux}. Se apoya en la misma \ac{API} REST documentada en el
anexo~\ref{anx:api-reference}, de modo que cualquier acción de la interfaz tiene su equivalente en un
\textit{endpoint}. Se organiza en cuatro pantallas: \textbf{Jobs}
envía grabaciones o transcripciones y gobierna su revisión hasta que se escriben en el grafo;
\textbf{Graph} navega, busca y explora la base de conocimiento, incluida la travesía de impacto de un
nodo; \textbf{Manual Actions} cubre las operaciones ajenas al pipeline —crear nodos y relaciones a
mano, anular contenido o forzar la reingesta de un trabajo—; y el \textbf{Admin Panel}, reservado a
la clave raíz, provisiona las claves de \ac{API} y sus roles.

Dos guías de usuario, mantenidas junto al código para que evolucionen con él, cubren el detalle de
cada superficie. El recorrido completo de la interfaz, con un escenario de ejemplo de principio a
fin, está en
\href{https://github.com/josepmafe/seshat/blob/main/docs/user-guide.md}{\texttt{docs/user-guide.md}};
la operación de los arneses de evaluación —ejecutarlos, leer los resultados en MLflow y calibrar los
umbrales— se recoge en
\href{https://github.com/josepmafe/seshat/blob/main/docs/eval-user-guide.md}{\texttt{docs/eval-user-guide.md}}.
